\documentclass[11pt,a4paper]{article}
\usepackage[T1]{fontenc}
\usepackage{lmodern,amsmath,amssymb,amsthm,microtype}
\usepackage[margin=27mm]{geometry}
\usepackage[colorlinks=true,urlcolor=blue]{hyperref}
\setlength{\emergencystretch}{2em}
\newtheorem{theorem}{Theorem}[section]
\newtheorem{lemma}[theorem]{Lemma}
\newtheorem{proposition}[theorem]{Proposition}
\newcommand{\Q}{\mathbb Q}\newcommand{\Z}{\mathbb Z}
\DeclareMathOperator{\Tr}{Tr}
\title{A detector normalization for a signed moving-contact cover}
\author{}\date{17 November 2026}
\begin{document}\maketitle
\begin{abstract}
A signed one-square cover gives twelve prime positions of exactly
the common-translation form in Sun's theorem. A rational normalizer
preserves the native six-test detector, ordered completeness and
whole-character guards. Its extraction conclusion loses only the
numerator-prime support of the height parameter. No moving-contact
height inequality, and hence no eleven-variable definition, is proved.
\end{abstract}

\section{The rational functions}

Fix distinct positive integers $c_1,\ldots,c_6$; one may take
$c_i=i$. For nonzero $a,t\in\Q$ adjoin a finite root
\[
 z^2=a^2-t,
\]
and put
\begin{equation}\label{eq:functions}
 \begin{split}
 F(t,a)&=\prod_{i=1}^6(t+2c_i a+c_i^2),\\
 \delta&=t/a,\qquad w=F/(4\delta^8),\\
 W&=\prod_{i=1}^6(1+C_i^2/F),\qquad
 V=\Tr_{\Q(\sqrt2)/\Q}(W^{-1}),\qquad b=4w/V^2.
 \end{split}
\end{equation}
The six $C_i$ are the actual conditional detector tests, with all
their original point and difference-chart guards. Guard $a,t,F,W,V$;
if the cover is to exclude its diagonal branch, also guard
$a^2-t=z^2$. The last condition is a scalar guard and adds no free
coordinate. The functions in \eqref{eq:functions} do not depend on
the choice of sign of $z$.

In the proposed height application,
\[
 M=(X^2+B)^e+T^e,\qquad a=M^{d/2},\qquad d\text{ even},
\]
where $T=\Tr x(A)$ and $e$ is even. Thus $M,a$ are positive
rational functions of retained abscissas and are the same at every
finite point-algebra character.

\section{Odd valuations and extraction}

Let $S$ contain two, the original fixed exceptional primes, and the
prime divisors of the fixed $c_i$. Write $v=v_p$ for $p\notin S$.

\begin{lemma}\label{lem:valuation}
Assume $z\in\Q$ and $F(t,a)\ne0$. If $v(F)$ is odd, then:
\begin{enumerate}
\item if $v(a)<0$, one has $v(t)=v(a)$ and $v(\delta)=0$;
\item if $v(a)=0$, one has $v(F)>0$ and $v(\delta)\ge0$;
\item if $v(a)>0$, one has $v(z)=0$, $v(t)=0$, $v(F)>0$
      and $v(\delta)=-v(a)$.
\end{enumerate}
\end{lemma}
\begin{proof}
Suppose first $q=v(a)<0$. If $v(z)<q$, then all six factors
have valuation $2v(z)$, so $v(F)=12v(z)$ is even. If $v(z)>q$,
all have valuation $2q$, again giving an even total. Thus an odd
total requires $v(z)=q$. Put $r=v(t)$. The terms $2c_i a$ have
valuation $q$ and the constants $c_i^2$ have valuation zero.
If $r<q$, every factor has valuation $r$, giving $6r$; if $r>q$,
every factor has valuation $q$, giving $6q$. Both are even.
Consequently $r=q$, proving the first assertion.

If $v(a)=0$ and $v(z)<0$, the total is again $12v(z)$.
Otherwise $t,F$ are integral, so an odd $v(F)$ is positive and
$v(\delta)=v(t)\ge0$.

Finally suppose $v(a)>0$. If $v(z)<0$, the total is even as
before. If $v(z)>0$, all six factors are units because every $c_i$
is a unit. Hence $v(z)=0$; then $t$ is a unit and $F$ is integral,
giving the remaining assertions.
\end{proof}

\begin{theorem}\label{thm:extraction}
Suppose the pole equation ensures that $b$ has no negative odd
valuation outside $S$. At every $p\notin S$ for which $v_p(F)$ is
odd and $v_p(a)\le0$, some place above $p$ satisfies
\[
                        v(C_i)>0\qquad(1\le i\le6).
\]
Moreover every negative valuation of $F$ outside $S$ is even in a
source solution.
\end{theorem}
\begin{proof}
Put $\nu=v_p(F)$ and $k=v_p(\delta)$. At every place above $p$,
each factor of $W$ has valuation
\[
                  \min(0,2v(C_i)-\nu)\le0;
\]
the two summands cannot tie because $\nu$ is odd. Let $q\le0$ be
the largest of the local valuations of $W$. Then
$v_p(V)\ge-q\ge0$. Consequently the pole condition gives
\[
 1\le v_p(b)\le\nu-8k+2q,
 \qquad q>-(\nu-8k)/2.
\]
At a place attaining $q$, each factor valuation is at least their
sum $q$, and hence
\begin{equation}\label{eq:bound}
                         v(C_i)>\nu/4+2k.
\end{equation}
If $v_p(a)<0$, Lemma~\ref{lem:valuation} gives $k=0$ and
the displayed inequality for $b$ forces $\nu>0$. If $v_p(a)=0$,
the lemma gives $\nu>0$ and $k\ge0$. Thus \eqref{eq:bound} is
positive in either case. A negative odd $\nu$ could only occur
in the first case, where it has just been excluded.
\end{proof}

The restriction on numerator primes cannot be removed from this
local argument. For example take $p=17$, $a=17^2$ and $z=18$.
Then $v_{17}(F)=1$ and $v_{17}(\delta)=-2$. With the toy tests
$C_1=1$, $C_2=\cdots=C_6=0$, one has $W=1+1/F\in\Q$ and
$V=2/W$, so $v_{17}(b)=15>0$ although $C_1$ is a unit.
This example concerns the valuation mechanism, not an actual
elliptic source tuple.

If $a=M^{d/2}$, its positive prime support is exactly that of the
numerator of $M$. Thus the omitted radical has cost
\[
 \log\prod_{v_p(a)>0}p
 =\log\prod_{v_p(M)>0}p\le h(M),
\]
independent of the exponent $d/2$. For a nonintegral residual
$X-n$, the native formal comparison therefore gives the prospective
radical bound
\[
 \log\operatorname{rad}_{\rm odd,S}(F)
 \le\log|\operatorname{num}(X-n)|+h(M),
\]
where only positive odd valuations outside $S$ enter the radical.
The point is the omission cost $h(M)$; no lower bound for this
moving-contact radical is asserted here.

\section{The exact prime construction}

\begin{proposition}\label{prop:prime}
Fix $a>0$ rational. Under Sun's Theorem 4.1, the signed cover has
infinitely many choices with $z>a+\max c_i$, for which $F$ is a
product of twelve distinct selected primes times a rational square.
The primes have the same simultaneous conjugate reduction-order,
avoidance and denominator-divisibility properties as in the native
conditional source.
\end{proposition}
\begin{proof}
Choose an integer $D>0$ such that every $D^2(a+c_i)$ is integral.
The twelve numbers $\pm D^2(a+c_i)$ are distinct. Apply Sun's
theorem to these twelve integer coefficients. It supplies $u,h>0$
such that
\[
 p_i^- =u-D^2(a+c_i)h^2,\qquad
 p_i^+ =u+D^2(a+c_i)h^2
\]
are the required primes. Put $z=u/(D^2h^2)$ and choose the slope
lower bound so that $z>a+\max c_i$. Then
\[
 (a+c_i)^2-z^2=-\frac{p_i^-p_i^+}{D^4h^4},\qquad
 F=\frac{\prod_i p_i^-p_i^+}{D^{24}h^{24}}>0.
\]
Both prime positions have the same $+u$ coefficient. Their signs
therefore impose no opposite Gaussian congruences: at a prime
dividing the prescribed denominator, both positions have the same
residue. The six minus signs cancel in the product.

The fixed input $a$ is chosen before prime selection. The avoidance
list may thus include the numerator and denominator prime support
of $a$, of every $c_i(2a+c_i)$, and of the rational coefficient
differences. At a selected prime the corresponding factor has
valuation one, all other factors are units, and $\delta=t/a$ is
a unit. Every other possible odd valuation of $F$ is absent.
\end{proof}

\section{Ordered completeness and the finite census}

Fix an integer input $X=n\ne0$, take $A=[n]P$, and choose the
same provisional point neighborhoods and signed-branch avoidance
as in the original conditional preparation. All quantities $M,a$
are now fixed. At each fixed finite place, as $z$ tends to infinity,
\[
 F\sim z^{12},\qquad \delta\sim-z^2/a,\qquad
 w\sim\frac{a^8}{4z^4},\qquad W\longrightarrow1,\quad
 V\longrightarrow2.
\]
The leading term of $w$ is a rational square. Arbitrary denominator
divisibility in Proposition~\ref{prop:prime} gives all the original
finite preparation depths, including the $2$-adic square class and
valuation and the added preparation at seven for the intrinsic pole.
At real places, $z>a+\max c_i$ gives $F>0$, and sufficiently large
$z$ gives small positive $w$ and $V>1$ in the prepared neighborhoods.
Thus $0<7b<1$ at the actual character. The twelve compatible contact
orders permit the same final refinement $H=[h_0]P$, $U=[n]H$.

At every fresh selected prime $v_p(F)=v_p(w)=1$, the detector has
its original successful slot, and $V$ is a unit. Thus $v_p(b)=1$.
At all other primes the weight valuation is even, except for the
already prepared fixed places. Hence the actual $7b$ has square
reduced denominator and satisfies the intrinsic realization
hypotheses. The intrinsic two-parameter pole theorem supplies
its rational witnesses and the established Zariski-density guard
avoidance.

At every point-algebra character, $a,t,F$ have the same rational
values, $F>0$, all factors of $W$ are positive, and $V>0$.
The new cover root is nonzero since $z>a+\max c_i$; the scalar
guards $a,t,F$ are nonzero. All original point, difference, quotient
and denominator guards remain present. The pole radicand norms
are avoided exactly as in the intrinsic transfer theorem.

The prospective finite algebra has rank $64\cdot2\cdot2=256$:
three point blocks, one cover root and one pole root. The proposed
unrestricted coordinates are one input, six point coordinates, one
retained $t$, two pole parameters and the single encoder coordinate.
This is a census of eleven coordinates for the incomplete
specification. Soundness still requires a new moving-contact height
theorem, with the omitted numerator primes of $M$ included. The
fixed-contact height theorem cannot be cited in its place.

\section{Remaining height problem}

Writing $r=z-a$ and $s=z+a$ gives
\[
 F=\prod_i(r-c_i)(s+c_i),\qquad s-r=2a.
\]
An odd prime common to a factor from each group may cancel from
the parity radical. Such a prime divides $2a+c_i+c_j$.
These are input-dependent collision factors, with heights of order
$h(a)$, so the old fixed-contact estimate does not apply. The
perfect-power constraint $a=M^{d/2}$ and the signed one-square
height lemma do not by themselves give the required lower bound.
In addition, the branch $z=0$ makes $F$ a square and must not be
silently included in a proposed radical-height inequality.

\begin{thebibliography}{9}\raggedright
\bibitem{Sun}Z.-W.~Sun, \emph{48 unknowns for Hilbert's tenth problem
over the rationals}, arXiv:2610.20691v3, Theorem 4.1,
\url{https://arxiv.org/pdf/2610.20691}.
\bibitem{Intrinsic}\emph{Thirteen-variable conditional sources with
intrinsic congruent-number poles}, repository manuscript,
\nolinkurl{papers/pole_parameters/conditional_intrinsic_transfer_20261117_v57.tex}.
\end{thebibliography}
\end{document}
